\chapter{单自由度系统强迫振动——简谐激励}
\intro{本章研究单自由度系统在简谐外力作用下的稳态响应。从无阻尼系统的强迫振动入手，引出拍振和共振现象;然后深入分析有阻尼系统的幅频特性、相频特性和频率响应函数;接着讨论旋转不平衡和基础激励两类重要的工程问题;最后介绍机械阻抗法和Nyquist图等频域分析方法。}

%=====================================================================
\section{无阻尼系统在简谐力作用下的响应}
%=====================================================================

\subsection{运动方程}

考虑单自由度无阻尼系统受简谐外力 $F(t) = F_0\sin\omega t$ 作用：
\begin{myformula}
    \[
    m\ddot{x} + kx = F_0\sin\omega t
    \]
\end{myformula}

方程两边除以 $m$，并利用 $\omega_n^2 = k/m$：
\begin{myformula}
    \[
    \ddot{x} + \omega_n^2 x = \frac{F_0}{m}\sin\omega t
    \]
\end{myformula}

\begin{center}
\begin{tikzpicture}[scale=1.2]
    \fill[vibWall] (0-0.2,0-0.8) rectangle (0+0.2,0+0.8);
    \draw[decorate,decoration={coil,aspect=0.5,amplitude=5,segment length=7}] (0.5,0) -- (2.3,0) node[midway,above] {$k$};
    \draw[fill=gray!30] (2.3,-0.5) rectangle (3.3,0.5);
    \node at (2.8,0) {$m$};
    
    \draw[-Stealth, red, thick] (3.3,0) -- (4.5,0) node[midway,above] {$F_0\sin\omega t$};
    
    \draw[-Stealth] (3.3,0.7) -- (5.0,0.7) node[right] {$x$};
    \draw[thick] (3.3,0.5) -- (3.3,0.9);
    \draw[thick] (-0.2,-0.7) -- (5.0,-0.7);
\end{tikzpicture}
\end{center}

\subsection{通解 = 齐次解 + 特解}

根据线性微分方程理论，全解由齐次解（自由振动）和特解（稳态强迫振动）叠加而成。

\textbf{齐次解}（自由振动部分）：
\begin{myformula}
    \[
    x_h(t) = A\cos\omega_n t + B\sin\omega_n t
    \]
\end{myformula}

\textbf{特解}：由于激励是正弦函数，设特解具有相同形式：
\begin{myformula}
    \[
    x_p(t) = X\sin\omega t
    \]
\end{myformula}

代入方程求 $X$。注意 $x_p$ 与激励频率 $\omega$ 相同（线性系统特性），仅是振幅和相位可能不同（无阻尼时相位差为 0 或 $\pi$）。

代入：$\ddot{x}_p = -\omega^2 X\sin\omega t$，得
\begin{myformula}
    \[
    -m\omega^2 X\sin\omega t + kX\sin\omega t = F_0\sin\omega t
    \]
    \[
    (k - m\omega^2)X = F_0
    \]
\end{myformula}

因此：
\begin{myformula}
    \[
    X = \frac{F_0}{k - m\omega^2} = \frac{F_0/k}{1 - (\omega/\omega_n)^2}
    \]
\end{myformula}

定义 $X_0 = F_0/k$ 为\textbf{静变形}（Static Deflection），即静力 $F_0$ 作用下系统的位移;定义 $r = \omega/\omega_n$ 为\textbf{频率比}（Frequency Ratio），则：
\begin{myformula}
    \boxed{X = \frac{X_0}{1 - r^2}}
    \end{myformula}

\subsection{全解与初始条件}

系统的全解为：
\begin{myformula}
    \[
    x(t) = A\cos\omega_n t + B\sin\omega_n t + \frac{F_0/k}{1 - r^2}\sin\omega t
    \]
\end{myformula}

由初始条件 $x(0) = x_0$，$\dot{x}(0) = v_0$：
\begin{myformula}
    \[
    A = x_0,\qquad B = \frac{v_0}{\omega_n} - \frac{r}{1-r^2}\cdot\frac{F_0}{k}
    \]
    \[
    x(t) = x_0\cos\omega_n t + \frac{v_0}{\omega_n}\sin\omega_n t
           - \frac{F_0/k}{1-r^2}\cdot\frac{\omega}{\omega_n}\sin\omega_n t
           + \frac{F_0/k}{1-r^2}\sin\omega t
    \]
\end{myformula}

特别地，零初始条件下（$x_0 = v_0 = 0$）：
\begin{myformula}
    \[
    x(t) = \frac{F_0/k}{1-r^2}\left(\sin\omega t - r\sin\omega_n t\right)
    \]
\end{myformula}

\subsection{拍振现象}

当激励频率 $\omega$ 接近（但不等于）固有频率 $\omega_n$ 时，即 $\omega \approx \omega_n$，全解中出现两个频率非常接近的简谐运动的叠加，从而产生\textbf{拍振}（Beating）现象。

零初始条件下：
\begin{myformula}
    \[
    x(t) = \frac{F_0/k}{1-r^2}(\sin\omega t - r\sin\omega_n t)
    \]
\end{myformula}

当 $\omega \approx \omega_n$ 时，利用三角恒等式可写为：
\begin{myformula}
    \[
    x(t) \approx \frac{F_0/k}{2(1-r)}\left[\sin\omega t - \sin\omega_n t\right]
          \approx \frac{F_0/m}{2\omega_n\varepsilon}\cdot 2\sin\frac{\omega_n - \omega}{2}t\cdot\cos\frac{\omega_n + \omega}{2}t
    \]
\end{myformula}

其中 $\varepsilon = \omega_n - \omega$。$x(t)$ 可视为频率为 $(\omega_n+\omega)/2$ 的振荡，其振幅按 $\sin[(\omega_n-\omega)t/2]$ 缓慢变化，形成"拍"。

拍的周期（振幅变化的周期）为：
\begin{myformula}
    \[
    T_b = \frac{2\pi}{|\omega_n - \omega|}
    \]
\end{myformula}

\textbf{物理意义}：拍振是两个频率接近的波叠加产生的干涉现象。在工程中，可以通过拍振来检测两个频率的微小差异——例如，可利用拍振方法对未知频率进行精确测量。

\subsection{共振}

当 $\omega = \omega_n$（即 $r = 1$）时，特解 $X = F_0/(k - m\omega^2)$ 的分母为零，特解形式不再适用。此时需要重新求特解。

\textbf{共振时的方程}：
\begin{myformula}
    \[
    \ddot{x} + \omega_n^2 x = \frac{F_0}{m}\sin\omega_n t
    \]
\end{myformula}

设特解形式为 $x_p(t) = X t\cos\omega_n t$（由于齐次解包含了 $\sin\omega_n t$ 和 $\cos\omega_n t$，特解需乘以 $t$ 以保证独立），代入原方程确定 $X$：
\begin{myformula}
    \[
    X = -\frac{F_0}{2m\omega_n}
    \]
\end{myformula}

因此共振特解为：
\begin{myformula}
    \boxed{x_p(t) = -\frac{F_0 t}{2m\omega_n}\cos\omega_n t}
    \end{myformula}

\textbf{共振的特征}：
\begin{itemize}
    \item 振幅随时间的增加而线性增长（无阻尼时）：$|x_p(t)| \propto t$
    \item 理论上当 $t \to \infty$ 时振幅趋于无穷大
    \item 实际系统中阻尼会限制振幅的增长，使其趋于有限值
    \item 无论从安全还是从使用角度看，共振通常是需要避免的（但也有利用共振的场合，如音叉、压电谐振器等）
\end{itemize}

%=====================================================================
\section{有阻尼系统在简谐力作用下的响应}
%=====================================================================

\subsection{运动方程}

考虑粘性阻尼的单自由度系统受简谐外力 $F(t) = F_0\sin\omega t$ 作用：
\begin{myformula}
    \[
    m\ddot{x} + c\dot{x} + kx = F_0\sin\omega t
    \]
\end{myformula}

标准形式：
\begin{myformula}
    \[
    \ddot{x} + 2\zeta\omega_n\dot{x} + \omega_n^2 x = \frac{F_0}{m}\sin\omega t
    \]
\end{myformula}

全解仍由齐次解（自由振动）和特解（稳态强迫振动）两部分组成。由于阻尼的存在，齐次解 $x_h(t)$ 是衰减的（$t \to \infty$ 时趋于零），因此在足够长时间后，只剩下特解 $x_p(t)$——称为\textbf{稳态响应}（Steady-State Response）。

\subsection{稳态解的求解}

设稳态解（特解）具有以下形式：
\begin{myformula}
    \[
    x_p(t) = X\sin(\omega t - \phi)
    \]
\end{myformula}
其中 $X$ 为稳态振幅，$\phi$ 为响应落后于激励的相位角（$0 \le \phi \le \pi$）。

\textbf{方法一：代入法}。将 $x_p$ 及其导数代入运动方程：
\[
\ddot{x}_p = -\omega^2 X\sin(\omega t-\phi)
\]
\[
\dot{x}_p = \omega X\cos(\omega t-\phi)
\]

代入 $m\ddot{x} + c\dot{x} + kx = F_0\sin\omega t$，展开三角函数并利用 $\sin(\omega t-\phi)$ 和 $\cos(\omega t-\phi)$ 的系数分别匹配，得到两个方程：
\begin{myformula}
    \[
    \begin{cases}
        (k - m\omega^2)X\cos\phi + c\omega X\sin\phi = F_0\\
        (k - m\omega^2)X\sin\phi - c\omega X\cos\phi = 0
    \end{cases}
    \]
\end{myformula}

\textbf{方法二：复数法}（更为简洁，见下文）。此处用代入法继续求解。

由第二个方程：
\begin{myformula}
    \[
    \tan\phi = \frac{c\omega}{k - m\omega^2} = \frac{2\zeta\omega/\omega_n}{1 - (\omega/\omega_n)^2} = \frac{2\zeta r}{1 - r^2}
    \]
\end{myformula}

由两个方程联立可得振幅：
\begin{myformula}
    \[
    X = \frac{F_0}{\sqrt{(k - m\omega^2)^2 + (c\omega)^2}}
    \]
\end{myformula}

整理为标准形式：
\begin{myformula}
    \boxed{X = \frac{F_0/k}{\sqrt{[1 - r^2]^2 + [2\zeta r]^2}}}
    \end{myformula}

\begin{myformula}
    \boxed{\phi = \arctan\frac{2\zeta r}{1 - r^2}}
    \end{myformula}

其中 $r = \omega/\omega_n$ 为频率比，$X_0 = F_0/k$ 为静变形。

\subsection{放大因子}

定义\textbf{放大因子}（Magnification Factor）或\textbf{动力放大系数} $M$：
\begin{myformula}
    \boxed{M = \frac{X}{X_0} = \frac{1}{\sqrt{[1 - r^2]^2 + [2\zeta r]^2}}}
    \end{myformula}

$M$ 的物理意义：稳态振幅与静力 $F_0$ 产生的静变形之比。$M$ 大于 1 表示动力放大效应，小于 1 表示隔振效果。

放大因子随频率比 $r$ 和阻尼比 $\zeta$ 的变化：
\begin{itemize}
    \item $r \to 0$（低频）：$M \to 1$，系统反应近似静力作用（刚度主导）
    \item $r \to 1$（共振区）：$M \to 1/(2\zeta)$（大振幅，阻尼决定峰值）
    \item $r \to \infty$（高频）：$M \to 0$，系统跟不上激励（惯性主导）
    \item 增大阻尼 $\zeta$ 可以显著降低共振区的放大因子
\end{itemize}

\subsection{幅频特性与相频特性}

\textbf{幅频特性} $M(r)$ 和\textbf{相频特性} $\phi(r)$ 是系统频率响应的两种表示，统称\textbf{频率响应特性}（Frequency Response Characteristics）。

\begin{center}
\begin{tikzpicture}[scale=1.0]
    % ===== 幅频特性 =====
    \begin{scope}
        \draw[-Stealth] (0,0) -- (0,4.5) node[above] {$M$};
        \draw[-Stealth] (0,0) -- (5.5,0) node[right] {$r = \omega/\omega_n$};
        \node[below left] at (0,0) {$0$};
        
        \draw[dashed, gray] (1,0) -- (1,4.3);
        \node[gray] at (1,-0.3) {$1$};
        
        % zeta = 0
        \draw[blue, very thick, domain=0:0.95, samples=100] plot (\x, {1/(1-\x*\x)});
        \draw[blue, very thick, domain=1.05:2.0, samples=100] plot (\x, {-1/(1-\x*\x)});
        
        % zeta = 0.1
        \draw[red, thick, domain=0:2.0, samples=100] plot (\x, {1/sqrt((1-\x^2)^2+(2*0.1*\x)^2)});
        
        % zeta = 0.3
        \draw[cyan!60!black, thick, domain=0:2.0, samples=100] plot (\x, {1/sqrt((1-\x^2)^2+(2*0.3*\x)^2)});
        
        % zeta = 0.7
        \draw[green!50!black, thick, domain=0:2.0, samples=100] plot (\x, {1/sqrt((1-\x^2)^2+(2*0.7*\x)^2)});
        
        % zeta = 1.0
        \draw[purple, thick, domain=0:2.0, samples=100] plot (\x, {1/sqrt((1-\x^2)^2+(2*1.0*\x)^2)});
        
        \node[blue] at (0.5,3.5) {$\zeta=0$};
        \node[red] at (1.5,2.0) {$\zeta=0.1$};
        \node[cyan!60!black] at (1.8,1.2) {$\zeta=0.3$};
        \node[green!50!black] at (2.0,0.8) {$\zeta=0.7$};
        \node[purple] at (2.2,0.5) {$\zeta=1.0$};
        
        \node at (2.8,4.2) {(a) 幅频特性};
    \end{scope}
    
    % ===== 相频特性 =====
    \begin{scope}[xshift=6.5cm]
        \draw[-Stealth] (0,0) -- (0,4.5) node[above] {$\phi$};
        \draw[-Stealth] (0,2.25) -- (5.5,2.25) node[right] {$r$};
        \node[below left] at (0,2.25) {$0$};
        
        \draw[dashed, gray] (1,0) -- (1,4.5);
        \node[gray] at (1,2.0) {$1$};
        
        \node at (0,4.5) {$\pi$};
        \node at (0,2.25) {$\pi/2$};
        
        % zeta = 0
        \draw[blue, very thick] (0,2.25) -- (0.98,2.25);
        \draw[blue, very thick] (1.02,4.5) -- (2.0,4.5);
        \fill[blue] (1,2.25) circle (0.06);
        \fill[blue] (1,4.5) circle (0.06);
        
        % zeta = 0.1
        \draw[red, thick, domain=0:2.0, samples=100] plot (\x, {2.25+rad(atan(2*0.1*\x/(1-\x*\x)))});
        
        % zeta = 0.3
        \draw[cyan!60!black, thick, domain=0:2.0, samples=100] plot (\x, {2.25+rad(atan(2*0.3*\x/(1-\x*\x)))});
        
        % zeta = 0.7
        \draw[green!50!black, thick, domain=0:2.0, samples=100] plot (\x, {2.25+rad(atan(2*0.7*\x/(1-\x*\x)))});
        
        % zeta = 1.0
        \draw[purple, thick, domain=0:2.0, samples=100] plot (\x, {2.25+rad(atan(2*1.0*\x/(1-\x*\x)))});
        
        \node[blue] at (0.5,1.8) {$\zeta=0$};
        \node[red] at (1.8,3.2) {$\zeta=0.1$};
        \node[cyan!60!black] at (2.0,4.0) {$\zeta=0.3$};
        \node[green!50!black] at (2.2,4.2) {$\zeta=0.7$};
        \node[purple] at (2.2,4.3) {$\zeta=1.0$};
        
        \node at (2.8,4.2) {(b) 相频特性};
    \end{scope}
\end{tikzpicture}
\end{center}

从相频特性可以看出：
\begin{itemize}
    \item $r \to 0$：$\phi \to 0$（响应与激励同相）
    \item $r = 1$：$\phi = \pi/2$（响应滞后激励 $90^\circ$），与 $\zeta$ 无关
    \item $r \to \infty$：$\phi \to \pi$（响应与激励反相）
    \item $\zeta$ 越大，相位变化越平缓
\end{itemize}

\subsection{共振频率与共振振幅}

放大因子 $M(r)$ 的最大值对应的频率称为\textbf{共振频率}（Resonance Frequency）。令 $\partial M/\partial r = 0$：
\begin{myformula}
    \[
    \frac{\partial M}{\partial r} = -\frac{1}{2}\cdot\frac{2[1-r^2](-2r) + 2[2\zeta r](2\zeta)}{\left([1-r^2]^2+[2\zeta r]^2\right)^{3/2}} = 0
    \]
\end{myformula}

即 $(1-r^2)(-r) + 2\zeta^2 r = 0$，得：
\begin{myformula}
    \boxed{r_{\text{res}} = \sqrt{1 - 2\zeta^2}}
    \end{myformula}

共振频率：
\begin{myformula}
    \[
    \omega_r = \omega_n\sqrt{1-2\zeta^2}
    \]
\end{myformula}

注意：
\begin{itemize}
    \item 有阻尼共振频率 $\omega_r$ 小于固有频率 $\omega_n$ 和有阻尼固有频率 $\omega_d = \omega_n\sqrt{1-\zeta^2}$
    \item 当 $\zeta > 1/\sqrt{2} \approx 0.707$ 时，$r_{\text{res}}$ 为虚数，$M(r)$ 无极值点，不出现共振峰
    \item 工程中阻尼比通常较小（$\zeta < 0.1$），共振频率与固有频率非常接近
\end{itemize}

\textbf{共振振幅}：将 $r_{\text{res}}$ 代入 $M(r)$：
\begin{myformula}
    \[
    M_{\max} = \frac{1}{\sqrt{[1 - (1-2\zeta^2)]^2 + [2\zeta\sqrt{1-2\zeta^2}]^2}}
             = \frac{1}{\sqrt{(2\zeta^2)^2 + 4\zeta^2(1-2\zeta^2)}}
    \]
    \[
    \boxed{M_{\max} = \frac{1}{2\zeta\sqrt{1-\zeta^2}}}
    \]
\end{myformula}

当 $\zeta \ll 1$ 时，近似有：
\begin{myformula}
    \[
    M_{\max} \approx \frac{1}{2\zeta}
    \]
\end{myformula}

\subsection{品质因数与半功率带宽}

\textbf{品质因数}（Quality Factor）$Q$ 是衡量系统共振锐度的参数：
\begin{myformula}
    \[
    Q = \frac{1}{2\zeta}
    \]
\end{myformula}

当 $\zeta$ 较小时，$Q \approx M_{\max}$。

\textbf{半功率带宽}（Half-Power Bandwidth）：在 $M(r)$ 曲线上，当 $M = M_{\max}/\sqrt{2}$ 时对应的两个频率 $r_1$ 和 $r_2$ 之间的宽度。当 $\zeta$ 较小时：
\begin{myformula}
    \[
    \Delta r = r_2 - r_1 \approx 2\zeta
    \]
\end{myformula}

由此可得\textbf{半功率带宽法测定阻尼比}：
\begin{myformula}
    \[
    \zeta \approx \frac{\Delta r}{2} = \frac{\omega_2 - \omega_1}{2\omega_n}
    \]
\end{myformula}

%=====================================================================
\section{频率响应函数（FRF）}
%=====================================================================

\subsection{机械阻抗}

\textbf{机械阻抗}（Mechanical Impedance）是频域中描述系统动态特性的重要概念。机械阻抗 $Z(\omega)$ 定义为激励力与响应速度的复数比：
\begin{myformula}
    \[
    Z(\omega) = \frac{F(\omega)}{\dot{X}(\omega)}
    \]
\end{myformula}

对于单自由度系统，将 $F = m\ddot{x} + c\dot{x} + kx$ 变换到频域（$F \to F e^{i\omega t}$, $x \to X e^{i\omega t}$）：
\begin{myformula}
    \[
    Z(\omega) = k + i\omega c - \omega^2 m
              = k\left[1 - r^2 + i\cdot 2\zeta r\right]
    \]
\end{myformula}

机械阻抗的三种分量：
\begin{itemize}
    \item \textbf{刚度阻抗}：$k$（与位移相关）
    \item \textbf{阻尼阻抗}：$i\omega c$（与速度相关）
    \item \textbf{质量阻抗}：$-\omega^2 m$（与加速度相关）
\end{itemize}

机械阻抗的倒数称为\textbf{机械导纳}（Mechanical Admittance）或\textbf{动柔度}（Receptance）。

\subsection{频率响应函数（位移导纳）}

定义\textbf{频率响应函数}（Frequency Response Function, FRF）$H(\omega)$ 为稳态位移响应与激励力的复数比：
\begin{myformula}
    \[
    H(\omega) = \frac{X(\omega)}{F(\omega)} = \frac{1}{k - m\omega^2 + i\omega c}
    \]
\end{myformula}

标准形式：
\begin{myformula}
    \boxed{H(\omega) = \frac{1/k}{1 - r^2 + i\cdot 2\zeta r}}
    \end{myformula}

$H(\omega)$ 的模即为放大因子/刚度：$|H(\omega)| = M/k$，$H(\omega)$ 的相角即为 $-\phi$。

FRF 包含了系统全部动力学特性信息：
\begin{itemize}
    \item \textbf{幅频} $|H(\omega)|$：反映不同频率下的响应大小
    \item \textbf{相频} $\angle H(\omega)$：反映响应与激励的相位关系
    \item \textbf{实频} $\Re\{H(\omega)\}$ 和\textbf{虚频} $\Im\{H(\omega)\}$
\end{itemize}

\subsection{Nyquist 图}

将 $H(\omega)$ 的实部作为横坐标、虚部作为纵坐标，以频率 $\omega$ 为参变量绘制的曲线称为 \textbf{Nyquist 图}（Nyquist Plot），也称矢量图或 Argand 图。

对于单自由度系统：
\begin{myformula}
    \[
    \Re\{kH(r)\} = \frac{1 - r^2}{(1-r^2)^2 + (2\zeta r)^2},\qquad
    \Im\{kH(r)\} = \frac{-2\zeta r}{(1-r^2)^2 + (2\zeta r)^2}
    \]
\end{myformula}

\begin{center}
\begin{tikzpicture}[scale=1.2]
    \draw[-Stealth] (-1.5,0) -- (1.5,0) node[right] {$\Re\{kH\}$};
    \draw[-Stealth] (0,-2.0) -- (0,0.5) node[above] {$\Im\{kH\}$};
    \node[below left] at (0,0) {$O$};
    
    % Nyquist circle for zeta=0.2
    \draw[blue, thick, domain=0:10, samples=200, variable=\x] 
        plot ({(1 - \x^2)/((1-\x^2)^2 + (2*0.2*\x)^2)}, {(-2*0.2*\x)/((1-\x^2)^2 + (2*0.2*\x)^2)});
    
    % 标注特殊点
    \fill[red] (1,0) circle (0.05) node[below right] {$r=0$};
    \fill[red] (0,-2.5) circle (0.05) node[above left] {$r=1$};
    
    \draw[-Stealth] (0.9,0) arc (0:-90:0.9);
    \node at (0.7,-0.5) {$\omega\uparrow$};
    
    \node at (1.0,-1.5) {$\zeta=0.2$};
    
    % 第二个圈
    \draw[red, thick, domain=0:10, samples=200, variable=\x] 
        plot ({0.7+(1 - \x^2)/((1-\x^2)^2 + (2*0.4*\x)^2)}, {0.7+(-2*0.4*\x)/((1-\x^2)^2 + (2*0.4*\x)^2)});
    \node[red] at (0.8,-0.2) {$\zeta=0.4$};
\end{tikzpicture}
\end{center}

Nyquist 图的特点：
\begin{itemize}
    \item 对于结构阻尼系统，Nyquist 图为正圆
    \item 对于粘性阻尼系统，Nyquist 图近似为圆（$\zeta$ 较小时）
    \item 可用于识别系统的模态参数（固有频率、阻尼比）
    \item 在实验模态分析中有广泛应用
\end{itemize}

%=====================================================================
\section{旋转不平衡引起的振动}
%=====================================================================

\subsection{等效力学模型}

旋转机械（电动机、涡轮机、压缩机等）由于转子质量偏心导致的不平衡力是工程中最常见的简谐激励源之一。

\begin{center}
\begin{tikzpicture}[scale=1.2]
    % 底座和弹簧-阻尼支撑
    \draw[thick] (-2,-0.5) -- (2,-0.5);
    
    \draw[decorate,decoration={coil,aspect=0.5,amplitude=4,segment length=6}] (-2,-0.5) -- (-2,0.5);
    \draw[decorate,decoration={coil,aspect=0.5,amplitude=4,segment length=6}] (2,-0.5) -- (2,0.5);
    \node at (-2.8,0) {$k$};
    \node at (2.8,0) {$k$};
    
    \draw (-0.5,0) rectangle (0.5,0.8) node[below] {-- (0,0)};
    
    % 机器质量
    \draw[fill=gray!30] (-1.2,0) rectangle (1.2,1.2);
    \node at (0,0.6) {$M$};
    
    % 旋转转子
    \draw[fill=gray!60] (0,0.6) circle (0.3);
    
    % 偏心质量
    \draw[fill=red] (0.2,0.95) circle (0.08) node[right] {$m$};
    
    \draw[Stealth-Stealth] (0,0.6) -- (0.2,0.95);
    \node at (0.0,0.8) {$e$};
    
    % 旋转箭头
    \draw[-Stealth] (0.3,0.7) arc (10:250:0.3);
    \node at (0.6,0.4) {$\omega t$};
    
    \draw[-Stealth] (1.5,0.6) -- (2.5,0.6) node[right] {$x$};
\end{tikzpicture}
\end{center}

设机器总质量为 $M$（不含偏心质量），偏心质量为 $m$，偏心距为 $e$，转子角速度为 $\omega$。系统的运动方程为（仅考虑铅垂方向）：
\begin{myformula}
    \[
    M\ddot{x} + c\dot{x} + kx = me\omega^2\sin\omega t
    \]
\end{myformula}

\textbf{注意}：右端的激振力幅值 $me\omega^2$ 正比于 $\omega^2$——这意味着在高转速下，即使偏心量很小，也会产生巨大的不平衡力。

将方程标准化：
\begin{myformula}
    \[
    \ddot{x} + 2\zeta\omega_n\dot{x} + \omega_n^2 x = \frac{me}{M}\omega^2\sin\omega t
    \]
\end{myformula}

稳态解 $x = X\sin(\omega t - \phi)$ 的振幅为：
\begin{myformula}
    \[
    X = \frac{me\omega^2/M}{\sqrt{(\omega_n^2 - \omega^2)^2 + (2\zeta\omega_n\omega)^2}}
    \]
\end{myformula}

无量纲化：
\begin{myformula}
    \boxed{\frac{MX}{me} = \frac{r^2}{\sqrt{(1-r^2)^2 + (2\zeta r)^2}}}
    \end{myformula}

相位角：
\begin{myformula}
    \[
    \phi = \arctan\frac{2\zeta r}{1-r^2}
    \]
\end{myformula}

\subsection{传递率}

\textbf{力传递率}（Force Transmissibility）$T_f$ 定义为传给基础的力的幅值与激振力幅值之比：
\begin{myformula}
    \[
    T_f = \frac{F_T}{F_0} = \frac{\sqrt{1 + (2\zeta r)^2}}{\sqrt{(1-r^2)^2 + (2\zeta r)^2}}
    \]
\end{myformula}

力传递率的特性：
\begin{itemize}
    \item $r < \sqrt{2}$：$T_f > 1$（力放大区）
    \item $r = \sqrt{2}$：$T_f = 1$（与阻尼无关）
    \item $r > \sqrt{2}$：$T_f < 1$（隔振区），且阻尼越小隔振效果越好
\end{itemize}

这是\textbf{隔振设计}的核心依据——要使隔振有效，必须使激励频率超过系统固有频率的 $\sqrt{2}$ 倍。

%=====================================================================
\section{基础激励（支座运动）}
%=====================================================================

\subsection{运动方程}

很多工程场合的振动不是由直接力引起的，而是由支座或基础的振动通过弹性元件传给系统。例如：建筑物在地震时的振动、车辆因路面不平引起的振动、精密仪器受环境振动的影响等。

\begin{center}
\begin{tikzpicture}[scale=1.2]
    % 基础
    \draw[thick] (-0.5,-1.0) -- (4.0,-1.0);
    \draw[-Stealth] (-0.5,-1.0) -- (-0.5,-1.8) node[below] {$y(t)=Y\sin\omega t$};
    
    \draw[decorate,decoration={coil,aspect=0.5,amplitude=4,segment length=6}] (0.5,-1.0) -- (0.5,0);
    \node at (0, -0.5) {$k$};
    
    \draw (0.8,0.2) rectangle (1.2,-0.2);
    \node at (1.8,0) {$c$};
    \draw (0.8,0) -- (1.2,-0.2);
    \draw (0.8,0) -- (1.2,0.2);
    
    \draw[fill=gray!30] (1.2,-0.4) rectangle (2.2,0.4);
    \node at (1.7,0) {$m$};
    
    \draw[-Stealth] (2.2,0.6) -- (3.5,0.6) node[right] {$x(t)$};
    \draw[thick] (2.2,0.4) -- (2.2,0.8);
    
    \draw[Stealth-Stealth, red] (-0.5,-1.5) -- (1.7,-1.5) node[midway,below] {相对运动};
\end{tikzpicture}
\end{center}

设基础位移为 $y(t) = Y\sin\omega t$，质量位移为 $x(t)$。

考虑弹簧力和阻尼力均取决于相对位移和相对速度：
\begin{itemize}
    \item 弹簧力：$F_k = k(x - y)$
    \item 阻尼力：$F_c = c(\dot{x} - \dot{y})$
\end{itemize}

由 Newton 第二定律：
\begin{myformula}
    \[
    m\ddot{x} = -k(x-y) - c(\dot{x}-\dot{y})
    \]
\end{myformula}

整理得：
\begin{myformula}
    \boxed{m\ddot{x} + c\dot{x} + kx = ky + c\dot{y}}
    \end{myformula}

将 $y = Y\sin\omega t$ 代入：
\begin{myformula}
    \[
    m\ddot{x} + c\dot{x} + kx = kY\sin\omega t + c\omega Y\cos\omega t
                              = Y\sqrt{k^2 + (c\omega)^2}\,\sin(\omega t + \alpha)
    \]
\end{myformula}
其中 $\alpha = \arctan(c\omega/k) = \arctan(2\zeta r)$。

\subsection{位移传递率}

位移传递率（Displacement Transmissibility）$T_d$ 定义为质量位移幅值与基础位移幅值之比：
\begin{myformula}
    \boxed{T_d = \frac{X}{Y} = \frac{\sqrt{1 + (2\zeta r)^2}}{\sqrt{(1-r^2)^2 + (2\zeta r)^2}}}
    \end{myformula}

相位角：
\begin{myformula}
    \[
    \phi = \arctan\frac{2\zeta r^3}{1 - r^2 + (2\zeta r)^2}
    \]
\end{myformula}

注意：力传递率与位移传递率的表达式完全相同！这是因为两者有对偶关系。

\subsection{隔振设计的原则}

由传递率曲线可知：
\begin{enumerate}
    \item 要使传递率 $T < 1$，必须 $r > \sqrt{2}$
    \item 在 $r > \sqrt{2}$ 的区域内，传递率随 $r$ 增大而单调递减
    \item 在隔振区域，阻尼越小则传递率越小，隔振效果越好——但阻尼又是控制共振振幅所必需的，需要在两者之间折中
    \item 工程设计中通常取 $r \in [3, 5]$，此时传递率在 $0.04 \sim 0.10$ 之间
\end{enumerate}

\subsection{相对运动}

在隔振设计中，另一个重要的量是相对位移 $z = x - y$。由运动方程可导出相对运动方程：
\begin{myformula}
    \[
    m\ddot{z} + c\dot{z} + kz = -m\ddot{y} = m\omega^2 Y\sin\omega t
    \]
\end{myformula}

相对位移的幅值为：
\begin{myformula}
    \[
    Z = \frac{m\omega^2 Y}{\sqrt{(k-m\omega^2)^2 + (c\omega)^2}}
      = \frac{r^2 Y}{\sqrt{(1-r^2)^2 + (2\zeta r)^2}}
    \]
\end{myformula}

这实际上与旋转不平衡激励下的形式相同（用 $\ddot{y}$ 代替 $me\omega^2/m$），因为本质都是基础的加速度引起的惯性力激励。

%=====================================================================
\section{复数频率响应法（机械阻抗法）}
%=====================================================================

\subsection{复数形式的运动方程}

用复数表示简谐激励力 $F(t) = F_0 e^{i\omega t}$（物理力取实部），设响应为 $x(t) = Xe^{i\omega t}$（$X$ 为复数，包含振幅和相位信息）。

代入运动方程：
\begin{myformula}
    \[
    (-m\omega^2 + i\omega c + k)X = F_0
    \]
\end{myformula}

\subsection{频率响应函数 $H(\omega)$}

由 $H(\omega) = X/F_0$：
\begin{myformula}
    \boxed{H(\omega) = \frac{1}{k - m\omega^2 + i\omega c}
                     = \frac{1/k}{1 - r^2 + i\cdot 2\zeta r}}
    \end{myformula}

$|H(\omega)|$ 即为位移幅值与力幅值之比：$|H(\omega)| = X/F_0 = M/k$。

\subsection{不同形式的频率响应函数}

根据响应的类型（位移、速度、加速度），有不同的频率响应函数：

\begin{itemize}
    \item \textbf{位移导纳}（Receptance）：$H(\omega) = X/F$
    \item \textbf{速度导纳}（Mobility）：$Y(\omega) = \dot{X}/F = i\omega H(\omega)$
    \item \textbf{加速度导纳}（Accelerance / Inertance）：$A(\omega) = \ddot{X}/F = -\omega^2 H(\omega)$
\end{itemize}

相应的倒数称为阻抗：
\begin{itemize}
    \item \textbf{位移阻抗}（Dynamic Stiffness）：$1/H(\omega)$
    \item \textbf{速度阻抗}（Mechanical Impedance）：$1/Y(\omega)$
    \item \textbf{加速度阻抗}（Apparent Mass）：$1/A(\omega)$
\end{itemize}

\subsection{复数频率响应的优势}

使用复数频率响应函数的优点：
\begin{enumerate}
    \item \textbf{运算简便}：微分变为代数运算（乘以 $i\omega$）
    \item \textbf{同时处理振幅和相位}：模 $|H|$ 给振幅信息，相角 $\angle H$ 给相位信息
    \item \textbf{可用于多自由度系统}：$H(\omega)$ 推广为频率响应函数矩阵
    \item \textbf{与实验模态分析对应}：实验测得的 FRF 可直接与解析模型比较
\end{enumerate}

%=====================================================================
\section{小结与对比}
%=====================================================================

\subsection{四种典型激励模型的对比}

\begin{center}
\begin{tabular}{c|c|c}
\hline
\textbf{激励类型} & \textbf{运动方程右端} & \textbf{无量纲振幅} \\
\hline
直接简谐力 & $\dfrac{F_0}{m}\sin\omega t$ & $\dfrac{X}{F_0/k} = \dfrac{1}{\sqrt{(1-r^2)^2+(2\zeta r)^2}}$ \\
\hline
旋转不平衡 & $\dfrac{me}{M}\omega^2\sin\omega t$ & $\dfrac{MX}{me} = \dfrac{r^2}{\sqrt{(1-r^2)^2+(2\zeta r)^2}}$ \\
\hline
基础激励 & $\omega_n^2 Y\sin\omega t + 2\zeta\omega_n\omega Y\cos\omega t$ & $\dfrac{X}{Y} = \dfrac{\sqrt{1+(2\zeta r)^2}}{\sqrt{(1-r^2)^2+(2\zeta r)^2}}$ \\
\hline
基础加速度 & $\omega^2 Y\sin\omega t$ & $\dfrac{Z}{Y} = \dfrac{r^2}{\sqrt{(1-r^2)^2+(2\zeta r)^2}}$ \\
\hline
\end{tabular}
\end{center}

\subsection{主要结论}

\begin{enumerate}
    \item 系统的稳态响应频率等于激励频率（线性系统特性）
    \item 阻尼使共振振幅有限，并使共振频率略低于固有频率
    \item 品质因数 $Q \approx 1/(2\zeta)$ 衡量共振锐度
    \item 隔振需确保激励频率超过 $\sqrt{2}$ 倍固有频率
    \item 频率响应函数是系统全部的频域特性描述，在实验和数值分析中都有核心地位
\end{enumerate}

%=====================================================================
\section{Python 代码：幅频特性与相频特性}
%=====================================================================

\begin{mycode}{绘制幅频特性与相频特性曲线}
\begin{lstlisting}[language=Python]
import numpy as np
import matplotlib.pyplot as plt

omega_n = 10.0  # 固有频率 rad/s
zeta_values = [0.0, 0.1, 0.2, 0.3, 0.5, 0.7, 1.0]
r = np.linspace(0.01, 3.0, 1000)

colors = plt.cm.viridis(np.linspace(0, 1, len(zeta_values)))

fig, (ax1, ax2) = plt.subplots(1, 2, figsize=(14, 5))

# 幅频特性
for zeta, color in zip(zeta_values, colors):
    M = 1 / np.sqrt((1 - r**2)**2 + (2 * zeta * r)**2)
    
    # 无阻尼时在 r=1 处分段处理，避免除零
    if zeta == 0:
        M_valid = np.where(np.abs(r - 1) > 0.02, M, np.nan)
    else:
        M_valid = M
    
    ax1.plot(r, M_valid, color=color, linewidth=1.5,
             label=f'$\\zeta = {zeta}$')

    # 标记共振峰
    if zeta < 1/np.sqrt(2) and zeta > 0:
        r_res = np.sqrt(1 - 2 * zeta**2)
        M_max = 1 / (2 * zeta * np.sqrt(1 - zeta**2))
        ax1.plot(r_res, M_max, 'o', color=color, markersize=5)
        ax1.annotate(f'$M_{{\\max}}={M_max:.2f}$',
                     (r_res, M_max), textcoords="offset points",
                     xytext=(0,10), fontsize=8, color=color)

ax1.set_xlabel('频率比 $r = \\omega/\\omega_n$')
ax1.set_ylabel('放大因子 $M$')
ax1.set_title('幅频特性')
ax1.legend(fontsize=8)
ax1.grid(True, alpha=0.3)
ax1.set_ylim(0, 6)

# 相频特性
for zeta, color in zip(zeta_values, colors):
    if zeta == 0:
        phi = np.where(r < 1, 0, np.pi)
    else:
        phi = np.arctan2(2 * zeta * r, 1 - r**2)
        phi = np.where(phi < 0, phi + np.pi, phi)
    
    ax2.plot(r, phi, color=color, linewidth=1.5,
             label=f'$\\zeta = {zeta}$')

ax2.set_xlabel('频率比 $r = \\omega/\\omega_n$')
ax2.set_ylabel('相位角 $\\phi$ (rad)')
ax2.set_title('相频特性')
ax2.legend(fontsize=8)
ax2.grid(True, alpha=0.3)
ax2.axhline(y=np.pi/2, color='gray', linestyle='--', alpha=0.5)
ax2.axhline(y=np.pi, color='gray', linestyle='--', alpha=0.5)
ax2.set_yticks([0, np.pi/4, np.pi/2, 3*np.pi/4, np.pi])
ax2.set_yticklabels(['$0$', '$\\pi/4$', '$\\pi/2$', '$3\\pi/4$', '$\\pi$'])

plt.tight_layout()
plt.show()

# 打印共振信息
print("=== 共振分析 ===")
print(f"固有频率: omega_n = {omega_n:.2f} rad/s")
for zeta in [0.05, 0.1, 0.2]:
    if zeta < 1/np.sqrt(2):
        r_res = np.sqrt(1 - 2*zeta**2)
        M_max = 1/(2*zeta*np.sqrt(1-zeta**2))
        print(f"zeta={zeta:.2f}: r_res={r_res:.4f}, M_max={M_max:.4f}, Q={1/(2*zeta):.4f}")
\end{lstlisting}
\end{mycode}

\begin{mycode}{旋转不平衡和基础激励的幅频特性}
\begin{lstlisting}[language=Python]
import numpy as np
import matplotlib.pyplot as plt

r = np.linspace(0.01, 4.0, 1000)
zeta_values = [0.05, 0.1, 0.2, 0.5, 1.0]

fig, (ax1, ax2) = plt.subplots(1, 2, figsize=(14, 5))

for zeta in zeta_values:
    # 直接简谐力
    M = 1 / np.sqrt((1 - r**2)**2 + (2*zeta*r)**2)
    
    # 旋转不平衡
    M_rot = r**2 / np.sqrt((1 - r**2)**2 + (2*zeta*r)**2)
    
    # 力传递率 / 位移传递率
    T = np.sqrt(1 + (2*zeta*r)**2) / np.sqrt((1 - r**2)**2 + (2*zeta*r)**2)
    
    ax1.plot(r, M_rot, linewidth=1.5, label=f'$\\zeta={zeta}$')
    ax2.plot(r, T, linewidth=1.5, label=f'$\\zeta={zeta}$')

ax1.set_xlabel('频率比 $r$')
ax1.set_ylabel('$MX/(me)$')
ax1.set_title('旋转不平衡振幅比')
ax1.legend()
ax1.grid(True, alpha=0.3)
ax1.axhline(y=1, color='gray', linestyle='--', alpha=0.5)

ax2.set_xlabel('频率比 $r$')
ax2.set_ylabel('传递率 $T$')
ax2.set_title('力传递率 / 位移传递率')
ax2.legend()
ax2.grid(True, alpha=0.3)
ax2.axhline(y=1, color='gray', linestyle='--', alpha=0.5)
ax2.axvline(x=np.sqrt(2), color='red', linestyle='--', alpha=0.5,
            label='$r=\\sqrt{2}$')
ax2.legend(fontsize=8)

plt.tight_layout()
plt.show()
\end{lstlisting}
\end{mycode}
